\documentclass[10pt]{article}
\usepackage[margin=20mm]{geometry}
\usepackage{amsmath,amssymb,mathtools,fontspec,unicode-math,xurl,hyperref}
\IfFontExistsTF{DejaVu Serif}{\setmainfont{DejaVu Serif}}{\setmainfont{Latin Modern Roman}}
\setmathfont{latinmodern-math.otf}
\setlength{\parindent}{0pt}\setlength{\parskip}{6pt}\sloppy
\title{Initial verification of
Group actions on contractible 2-complexes II}
\author{}
\date{}
\begin{document}
\maketitle
Abstract: The local representation-theoretic calculations, central-extension argument, and two degree-zero arguments were checked. No substantive error was found within this scope. The report concerns arXiv:2102.11459v2, not an independently compared publisher PDF. External structural results remain identified dependencies.\par
\section*{Identity and review scope}
Paper M25-0101 is by Kevin Iván Piterman and Iván Sadofschi Costa, with journal DOI 10.1007/s00222-025-01363-8. The reviewed file is the seven-page arXiv version dated 21 August 2025 [paper]. Page numbers below are the PDF’s one-based pages, identical to its printed pages. All seven pages were read and visually inspected; page 7 contains affiliations rather than an argument.\par
Theorem C concerns finite acyclic, fixed-point-free, two-dimensional \(G\)-complexes, where \(G=\mathrm{PSL}_2(q)\), \(q>5\), and \(q\equiv5\) or \(13\pmod {24}\). Here \(q\) is a prime power, as required by the group notation. Its conclusion is the existence of a nontrivial unitary representation of the fundamental group. The passage to the fixed-point theorem for arbitrary finite groups uses Part I and the structural reductions cited there. The present review does not replace a review of that entire dependency chain.\par
\section*{Verified reasoning}
\subsection*{Simultaneous centralizer bound}
Lemma 2.1 (p. 2) does not assert that two arbitrary group elements commute. Instead it diagonalizes their images separately, using the same choice on isomorphic irreducible summands. In an isotypic decomposition, these diagonalizing operators act on representation factors and commute with the multiplicity-space operators. Thus \(A_1^{-1}A_2\) commutes with the original representation’s centralizer.\par
In the resulting common coordinate basis let \(r_{\lambda,\mu}\) count the coordinates with the indicated pair of eigenvalues. Then \(\sum r_{\lambda,\mu}=m\). The intersection of the two diagonal centralizers contains \(\prod_{\lambda,\mu}U(r_{\lambda,\mu})\), so \[\dim_{\mathbb R}(C_1\cap C_2)\geq\sum_{\lambda,\mu}r_{\lambda,\mu}^2
 \geq\frac{m^2}{n_1n_2}.\] The last inequality is Cauchy–Schwarz, with zero-sized blocks allowed. This verifies the bound and the compatibility condition later needed for the homotopy.\par
\subsection*{The character calculations omitted in Proposition 4.2}
Set \(m=(q-1)/2\) and \[a=\frac{-1+\sqrt q}{2},\qquad b=\frac{-1-\sqrt q}{2}.\] For the character listed in Proposition 4.1 (pp. 3–4), \(a^2+b^2=(q+1)/2\). Using precisely the class counts stated there gives \[\begin{aligned}
 \langle\chi|_B,\chi|_B\rangle_B
 &=\frac{2m^2+(q-1)(a^2+b^2)}{q(q-1)}=1,\\
 \langle\chi|_{C_4},\chi|_{C_4}\rangle_{C_4}
 &=\frac{2m^2}{4}=\frac{(q-1)^2}{8},\\
 \langle\chi|_{2D_{q+1}},\chi|_{2D_{q+1}}\rangle_{2D_{q+1}}
 &=\frac{2m^2+(q-1)}{2(q+1)}=\frac{q-1}{4}.\end{aligned}\] The \(C_4\) noncentral elements have character zero. The \(q+1\) order-four elements in the binary dihedral subgroup also contribute zero, while the remaining paired noncentral cyclic elements have character of absolute value one. For a unitary representation, the squared character norm is the sum of squared irreducible multiplicities, which is the real dimension of its unitary centralizer. The calculations therefore give all three conclusions of Proposition 4.2. They check the use of the supplied character table; they do not independently construct Dornhoff’s character table.\par
\subsection*{Central extension and differential}
For \(\widehat G=\mathrm{SL}_2(q)\), write \(z=-I\). In Lemma 6.1 (pp. 4–5), the homomorphism \(\mathbb C[\widehat G]\to\mathbb C[G]\) has kernel generated by \(1-z\). Every lifted edge stabilizer has twice the order of its image, hence \[\pi\bigl(N(\widehat G_e)\bigr)=2N(G_e).\] The factor \(1/2\) used to lift the coefficients is consequently required and correct. Applying the resulting group-ring identity to a representation gives the stated spanning relation.\par
For the adjoint representation on \(\mathfrak u(m)\), \(\rho_0(z)=-I\) acts trivially. The \((1-z)\) term disappears, giving the spanning of the tangent space by the images of the centralizer tangent spaces. Although the displayed group-ring identity uses complex coefficients, it may be applied to \(\mathfrak u(m)\otimes_{\mathbb R}\mathbb C\); full rank of this complexification implies full rank of the original real differential. Thus this change of coefficient field does not supply a counterexample.\par
The differential formula in Part I, Theorem 5.7, was checked in its versioned HTML [part1]. Its common minus sign does not affect surjectivity. The graph action is an admissible \(G\)-CW action; it is not an unrestricted action on a topological space.\par
\subsection*{Quotients and the degree-zero arguments}
In Proposition 7.1 (pp. 5–6), the three-factor product map is invariant under the displayed action: substituting the three transformed inverse factors cancels the \(\alpha_{v_1}\) and \(\alpha_{v_3}\) factors consecutively. The action is free, since the first and last fixed-coordinate equations successively force both acting elements to be the identity. Compactness makes it proper. The excess of the quotient dimension over the target dimension is \[\frac{q-1}{4}-\frac{(q-1)^2}{8}
 =\frac{(q-1)(3-q)}8<0
 \quad(q\equiv13\pmod {24}).\] A smooth map from this lower-dimensional compact manifold cannot be surjective, giving degree zero for the original map with the same image.\par
For Proposition 7.2 (p. 6), the cyclic edge groups have orders 6 and 4, so their numbers of eigenvalues are at most 6 and 4. Lemma 2.1 therefore gives \[\dim K\geq \frac{m^2}{24}=\frac{(q-1)^2}{96}.\] The relative conjugator commutes with the vertex centralizer at \(v_2\). Along the tree path \(\eta_0,\eta_2,\eta_3\), the other conjugators are identical. These are the compatibility conditions in Part I, Proposition 5.11; its connected-double-centralizer condition holds for unitary centralizers [part1].\par
For the enlarged action in Proposition 7.2, the fixed first, second and third coordinates force \(\alpha_1=\alpha_K=\alpha_3=I\), respectively. Also, \[(\alpha_1^{-1}\tau_0)^{-1}
 (\alpha_3^{-1}\tau_2\alpha_1)^{-1}
 (\alpha_K^{-1}\tau_3\alpha_3)^{-1}
 (\tau_1\alpha_K)^{-1}
 =\tau_0^{-1}\tau_2^{-1}\tau_3^{-1}\tau_1^{-1}.\] This verifies the quotient factorization without assuming commutativity of the factors. The dimension excess is bounded above by \[\frac{q-1}{4}-\frac{(q-1)^2}{96}
 =\frac{(q-1)(25-q)}{96}<0.\] Indeed, \(q\equiv5\pmod {24}\) and \(q>5\) imply \(q\geq29\). The excluded value \(q=5\) is not a counterexample to the stated proposition.\par
\section*{Dependence, outcome, and remaining work}
The dimension identity for the representation moduli is compatible with the lifted action: the two-cell stabilizer becomes \(\langle z\rangle\), whose scalar image has centralizer \(U(m)\). The Borel and full-group centralizer dimensions are both one, so their terms cancel in the equivariant Euler-character identity. Part I’s Corollary 8.2 and Proposition 8.3 were consulted for this step.\par
The checked arguments preserve the two ingredients used in Section 5: a regular point and a degree-zero map. The global reduction, existence of the acyclic attaching construction, and all results imported from Part I were not re-proved from first principles. No correction is proposed.\par
A bounded title/erratum search and the publisher and arXiv records were consulted on 26 October 2026 UTC. No separate correction was identified in those results. This is not an exhaustive absence claim. No comparison with the publisher’s eight-page typeset article was performed. No empirical experiment or author-code execution is relevant to the calculations above.\par
Item | Conclusion\par
Issue judgment | No Issue Found within the checked local proof chain.\par
Supported impact | No demonstrated damage to Theorem C.\par
Repair verification | Not applicable; no repair proposed.\par
Adoption | Not requested; scientific decisions remain pending.\par
Remaining work | Independent acceptance; full external dependency audit if required; comparison with the journal PDF.\par
\section*{Supporting checks}
All calculations above are analytic checks, not finite experiments or a formal proof-assistant certification. Source PDF pages 1–7 were actually rendered and inspected. A Poppler `Symbol` display-font warning was recorded; the inspected pages retained readable mathematics.\par
> f971f0f98871e3178aa9ee89eb1f530a5902cf0e506fac6d3ff8ac7a3bd53821.\par
Exact source-manifest, compiler, output and page-inspection bindings are in the accompanying provenance and coverage records. The current assistant performed the scientific checks. No human scientific approval is asserted.\par
9 K. I. Piterman and I. Sadofschi Costa, Group actions on contractible 2-complexes II, arXiv:2102.11459v2, 21 August 2025. https://arxiv.org/abs/2102.11459v2. I. Sadofschi Costa, Group actions on contractible 2-complexes I, arXiv:2102.11458v2, 21 August 2025. Selected dependency statements and proofs consulted in https://arxiv.org/html/2102.11458v2; not a full independent review of Part I. Springer Nature, publisher record for DOI 10.1007/s00222-025-01363-8, accessed 26 October 2026 UTC. Used for metadata and a bounded correction check.\par
\end{document}
